% Source: https://artofproblemsolving.com/community/c3h2153760_digit_sums
% Stable identity: aops-topic-2153760-q01
Write the original number as $N=5\cdot10^{n-1}+x$, where $0\le x<10^{n-1}$ is the number formed by the last $n-1$ digits. The new number is $10x+5$.

\textbf{(a)} Since $N\ge5\cdot10^{n-1}$, we get $3N\ge15\cdot10^{n-1}>10^n$, but the new number $10x+5$ is less than $10\cdot10^{n-1}=10^n$. So the new number is always less than $3N$.

\textbf{(b)} We need $3(10x+5)=N=5\cdot10^{n-1}+x$, that is,
\[29x=5\cdot10^{n-1}-15=5\left(10^{n-1}-3\right).\]
Since $\gcd(29,5)=1$, this has an integer solution exactly when $29\mid10^{n-1}-3$, and then $x=\dfrac{5\left(10^{n-1}-3\right)}{29}$ automatically satisfies $0\le x<10^{n-1}$.

Now $10\cdot3=30\equiv1\pmod{29}$, so $3\equiv10^{-1}$, and $10^{n-1}\equiv3$ is equivalent to $10^{n}\equiv1\pmod{29}$. The order of $10$ modulo the prime $29$ divides $28$. Checking the divisors:
\[10^1\equiv10,\qquad 10^2=100\equiv13,\qquad 10^4\equiv13^2=169\equiv24,\]
\[10^7=10^4\cdot10^2\cdot10\equiv24\cdot13\cdot10\equiv17,\qquad 10^{14}\equiv17^2=289\equiv28\pmod{29},\]
none of which is $1$, so the order is $28$. Hence the smallest $n$ is $\boxed{28}$, and the number is
\[N=5\cdot10^{27}+\dfrac{5\left(10^{27}-3\right)}{29}=\boxed{5172413793103448275862068965}.\]
Indeed, moving the $5$ to the end gives $1724137931034482758620689655$, and $3\cdot1724137931034482758620689655=5172413793103448275862068965$.
